\chapter{Ordered Taylor maps for supported currents}
\label{chap:ordered-current-taylor}

The quantum differential permits multiplication of observables with
disjoint supports.  A finite change of current representatives makes
their Wick products strict.  Interval interpolation then gives ordered
Taylor maps into the actual quantum complex.  Their linear term has
an explicit compact homotopy to the common-annulus current inclusion.

\section{The support of an ordered product}

Fix a centered disk $U=D_R\subset\C$, distinct points $a_1,\ldots,a_m\in U$,
and disjoint centered disks $D_i\Subset U$ about these points.
Choose smooth radial cutoffs $\chi_i\in C_c^\infty(D_i)$ equal to
one near $a_i$.  Choose an outer radial cutoff $\chi\in C_c^\infty(U)$
equal to one on a disk containing every $\overline D_i$.
For each $i$, fix a finite set $J_i\subset\Z_{\ge0}$ of jet orders.
Let $E=\{(i,n):1\le i\le m,\ n\in J_i\}$.
For $\nu=(i,n)\in E$, put
\[
 u_\nu=-\frac{\bar\partial\chi_i}{(z-a_i)^{n+1}},\qquad
 e_\nu=-\frac{\bar\partial\chi}{(z-a_i)^{n+1}},\qquad
 f_\nu=\frac{\chi-\chi_i}{(z-a_i)^{n+1}}.
\]
These are compact smooth representatives, and
$\bar\partial f_\nu=u_\nu-e_\nu$.
Their observable degrees are $|u_\nu|=|e_\nu|=0$ and $|f_\nu|=-1$.
Use the increasing orientation on $[0,1]$ and the complex orientation
on $U$, with residue integral $(2\pi i)^{-1}\int$.

Fix the bulk parameter $\ell\in\C$ and work first over
$R=\C[\hbar]$.  The boundary complex is
\[
 \mathcal B_\ell(V)
 =\left(\operatorname{Sym}(\Omega_c^{0,\bullet}(V)[1])[\hbar],
       D=\bar\partial+\hbar\Delta_{\mu_\ell}\right),\qquad
 \mu_\ell(v,w)=-\ell\int_V^{\mathrm{res}}v\wedge\partial w.
\]
The shifted Dolbeault differential has the sign specified in
Lemma~\ref{lem:bf-boundary-differential}.
Each symmetric power uses compact smooth kernels on $V^r$ and graded
coinvariants.  The sum over symmetric powers is algebraic.
No field-number completion is used.
Define symmetric constants, with a total order on $E$, by
\[
 p_{(i,n),(j,r)}=
 \begin{cases}
 \displaystyle\frac{\ell}{n!r!}
   \partial_{a_i}^{n}\partial_{a_j}^{r}(a_i-a_j)^{-2},&i\ne j,\\[3pt]
 0,&i=j.
 \end{cases}
\]
The extracted current level is $k=-\ell$.  Thus its divided-jet
Wick contraction is $-\hbar p_{\nu\omega}$.

Write $x_\nu,y_\nu,\theta_\nu$ for the generators represented by
$u_\nu,e_\nu,f_\nu$, respectively.  Their finite compact complex is
\[
 Q=R[x_\nu,y_\nu\mid\nu\in E]\otimes\Lambda(\theta_\nu\mid\nu\in E),
 \qquad |\theta_\nu|=-1,
\]
with left odd derivatives and differential
\begin{equation}
\label{eq:ordered-current-D}
 D_Q=\sum_\nu(x_\nu-y_\nu)\partial_{\theta_\nu}
 +\hbar\sum_{\nu<\omega}p_{\nu\omega}
 (\partial_{x_\omega}\partial_{\theta_\nu}
  +\partial_{x_\nu}\partial_{\theta_\omega}).
\end{equation}
The substitution $\rho:Q\hookrightarrow\mathcal B_\ell(U)$ is the
embedding of Proposition~\ref{prop:bf-compact-current-complex}.
For a graded $R$-module $M$, write
$T^c(M)=\bigoplus_{n\ge1}M^{\otimes_R n}$, with coproduct given
by splitting words between consecutive entries.

\begin{proposition}[The overlapping product obstruction]
\label{prop:ordered-current-overlap-obstruction}
Suppose $\hbar p_{\nu\omega}\ne0$ in $R$.
On $T^c(sQ)$ there is no uncurved square-zero coderivation whose
unary and binary terms are
\[
 B_1(sa)=sD_Qa,\qquad B_2(sa,sb)=(-1)^{|a|}s(ab),
 \qquad |sa|=|a|-1.
\]
The assertion holds regardless of the terms of arity at least three.
\end{proposition}
\begin{proof}
The component from length two to length one of the square is
\[
 (B_1B_2+B_2B_1)(sa,sb)
 =(-1)^{|a|}s\bigl(D_Q(ab)-(D_Qa)b-(-1)^{|a|}aD_Qb\bigr).
\]
Here the second unary insertion has sign $(-1)^{|sa|}$.
For $a=\theta_\nu$, $b=x_\omega$, equation
\eqref{eq:ordered-current-D} makes the right side
$-\hbar p_{\nu\omega}s1$.
It is nonzero.  An operation of arity at least three vanishes
on a word of length two.  Hence no such operation changes this square.
\end{proof}

For $S\subseteq\{1,\ldots,m\}$, define
$U_S=\coprod_{i\in S}D_i$ and $E_S=\{(i,n)\in E:i\in S\}$.
The binary factorization product
\[
 \mu_{S,T}:\mathcal B_\ell(U_S)\otimes_R\mathcal B_\ell(U_T)
       \longrightarrow\mathcal B_\ell(U_{S\cup T})
\]
exists for $S\cap T=\varnothing$.  It extends each input by zero
and multiplies its observable kernels.
Every cross contraction is zero: the local pairing evaluates two
generators at a common point, and their supports are disjoint.
Thus $D$ obeys the graded Leibniz rule for this product.
The same conclusion holds on compact smooth kernels by restriction
to the relevant diagonals.  These products are associative and
commute with the indicated extensions by zero.

\section{A compact lift of every Wick polynomial}

Let $W_S=R[y_\nu\mid\nu\in E_S]$, concentrated in degree zero.
The variables are normal symbols for divided holomorphic jets.
This is the finite-current specialization of the rational Wick
symbols of Chapter~\ref{chap:free-field-interval}, under
\[
 z_i\longmapsto a_i,\qquad \kappa\longmapsto-\ell,\qquad
 x_{i,n}\longmapsto n!y_{(i,n)}.
\]
Here $\kappa$ is the algebraic current level of that chapter.
The map evaluates rational coefficients before any radial Laurent
expansion.  Distinct $a_i$ make all denominators defined.
For disjoint $S,T$, set
\[
 K_{S,T}=\sum_{\nu\in E_S,\ \omega\in E_T}
       p_{\nu\omega}\partial_{y_\nu}\partial_{y_\omega},\qquad
 f\star g=e^{-\hbar K_{S,T}}(fg),
 \qquad
 K_S=\sum_{\substack{\nu<\omega\\\nu,\omega\in E_S}}
       p_{\nu\omega}\partial_{y_\nu}\partial_{y_\omega}.
\]
All derivatives fix the coefficients.  Each exponential terminates
on a polynomial.  The commuting operators satisfy
$K_{S\cup T}=K_S+K_T+K_{S,T}$.
It follows that $\star$ is associative on pairwise disjoint labels.

Let $\rho_S^x:W_S\to\mathcal B_\ell(U_S)$ substitute $u_\nu$ for
$y_\nu$ and multiply the resulting form generators.
Define the degree-zero linear map
\begin{equation}
\label{eq:ordered-current-Gamma}
 \Gamma_S(f)=\rho_S^x(e^{\hbar K_S}f).
\end{equation}

\begin{theorem}[Separated Wick lift]
\label{thm:ordered-current-wick-lift}
The maps $\Gamma_S$ are chain maps and satisfy
\begin{equation}
\label{eq:ordered-current-strict-product}
 \mu_{S,T}(\Gamma_Sf,\Gamma_Tg)=\Gamma_{S\cup T}(f\star g)
 \qquad(S\cap T=\varnothing).
\end{equation}
Their values lie in the actual polynomial current subcomplex
$R[u_\nu\mid\nu\in E_S]$ of $\mathcal B_\ell(U_S)$.
They are linear isomorphisms onto that subcomplex.
For $S\subseteq T$, let $\jmath_{S,T}$ extend observables by zero
from $U_S$ to $U_T$, and include $W_S$ into $W_T$ by the same
normal variables.  Then $\jmath_{S,T}\Gamma_S=\Gamma_T|_{W_S}$.
\end{theorem}
\begin{proof}
The differential kills every polynomial in form generators: their
Dolbeault derivatives vanish, and $\mu_\ell$ pairs no two such generators.
This proves the chain-map assertion.
For the product, compute before substitution:
\[
 e^{\hbar K_{S\cup T}}e^{-\hbar K_{S,T}}(fg)
 =e^{\hbar(K_S+K_T)}(fg)
 =(e^{\hbar K_S}f)(e^{\hbar K_T}g).
\]
Substitution turns the last product into the disjoint factorization
product.  This proves \eqref{eq:ordered-current-strict-product}.
The $u_\nu$ are independent by their disjoint disks and angular
Fourier powers.  Their polynomial algebra embeds in smooth kernels,
as in Proposition~\ref{prop:bf-compact-current-complex}.
The inverse exponential $e^{-\hbar K_S}$ proves the isomorphism assertion.
On a polynomial using only $E_S$, every derivative in $K_T-K_S$
kills that polynomial and every successive derivative of it.
Thus $e^{\hbar K_T}|_{W_S}=e^{\hbar K_S}$, which proves the
extension-by-zero identity.
\end{proof}

For currents at two and three distinct points, with $u_i=u_{(i,0)}$,
\[
 \begin{aligned}
 \Gamma_{\{1,2\}}(y_1y_2)&=u_1u_2+\hbar p_{12},\\
 \Gamma_{\{1,2\}}(y_1\star y_2)&=u_1u_2,\\
 \Gamma_{\{1,2,3\}}(y_1\star y_2\star y_3)&=u_1u_2u_3.
 \end{aligned}
\]
The nonzero contractions appear when these supported products are
expressed in common-annulus coordinates.

\section{The ordered bar and its interval maps}

Let $A_I=\C[t,dt]$, with $|t|=0$, $|dt|=1$, $d_It=dt$,
$d_I(dt)=0$, and $(dt)^2=0$.  Let
$C_I=\C e_0\oplus\C e_1\oplus\C\eta$ have degrees $0,0,1$ and
differential $de_0=-\eta$, $de_1=\eta$, $d\eta=0$.
The contraction is
\[
 \begin{aligned}
 i(e_0)&=1-t,&i(e_1)&=t,&i(\eta)&=dt,\\
 p(f)&=f(0)e_0+f(1)e_1,&
 p(gdt)&=\left(\int_0^1g\right)\eta,&
 h(gdt)&=\int_0^t g(u)du-t\int_0^1g(u)du,
 \end{aligned}
\]
with $h(f)=0$.  Direct differentiation gives
$dh+hd=1-ip$ and $hi=ph=h^2=0$.
There are no endpoint vanishing conditions.
Set
\[
 C(S)=C_I\otimes W_S,\quad A(S)=A_I\otimes W_S,\quad
 \mathcal T(S)=A_I\otimes\mathcal B_\ell(U_S).
\]
The source differential acts on the interval factor.
On the actual target, use
\[
 d_{\mathcal T}(\alpha\otimes q)
   =d_I\alpha\otimes q+(-1)^{|\alpha|}\alpha\otimes Dq,
\]
and, for disjoint $S,T$, the product
\[
 (\alpha\otimes q)(\beta\otimes r)
 =(-1)^{|q||\beta|}(\alpha\wedge\beta)\otimes\mu_{S,T}(q,r).
\]
The total differential is a derivation for these disjoint products.
Define $G_S:A(S)\to\mathcal T(S)$ by
$G_S(\alpha\otimes f)=\alpha\otimes\Gamma_Sf$.
Theorem~\ref{thm:ordered-current-wick-lift} makes $G$ a strict
morphism of these systems of differential graded products.

For any such collection $V(S)$, its supported tensor coalgebra is
\[
 T^c_{\mathrm{disj}}(sV)=
 \bigoplus_{n\ge1}\ \bigoplus_{S_1,\ldots,S_n\ \mathrm{pairwise\ disjoint}}
 sV(S_1)\otimes_R\cdots\otimes_R sV(S_n).
\]
The coproduct cuts a word between consecutive entries.  Empty label
sets may occur and carry the unit.  Each word has finite length.
A Taylor term replaces a consecutive block by its union label.
It preserves disjointness from all remaining entries.

For a system with differential $d$ and product $\mu$, define
$B_1(sa)=sda$ and $B_2(sa,sb)=(-1)^{|a|}s\mu(a,b)$.
Its ordered bar differential on $u_1|\cdots|u_n$ is
\begin{equation}
\label{eq:ordered-current-bar}
 B(u_1|\cdots|u_n)=
 \sum_{q=1}^{2}\sum_{r=0}^{n-q}
 (-1)^{\sum_{j=1}^{r}|u_j|}
 u_1|\cdots|u_r|B_q(u_{r+1},\ldots,u_{r+q})|
 u_{r+q+1}|\cdots|u_n.
\end{equation}
The degree is one.  The differential, Leibniz, and associativity
identities give $B^2=0$, in lengths one, two, and three, respectively.
There are no other components of the square.
This is the suspended convention for the tensor bar construction
described by Keller~\cite{KellerAInfinity}, \S3.6, pp.~11--12.

Write $I,P,H$ for the suspended interval maps on $sA$.
Let $B^A$ and $B^{\mathcal T}$ be the bar differentials on $A$ and
$\mathcal T$.  On $C$ put $b_1=s d_C s^{-1}$ and define recursively
\[
 \begin{split}
 f_1&=I,\qquad
 R_n=\sum_{a+b=n}B^A_2(f_a\otimes f_b),\\
 f_n&=-HR_n,\qquad b_n=PR_n\quad(n\ge2),\qquad
 \mathcal F_n=(sG s^{-1})f_n.
 \end{split}
\]
Each $f_n$ and $\mathcal F_n$ has suspended degree zero.

\begin{theorem}[Ordered current Taylor comparison]
\label{thm:ordered-current-taylor}
The maps $b_n$ define an uncurved $A_\infty$ structure on $C(S)$.
The maps $\mathcal F_n$ define a morphism of differential graded
supported tensor coalgebras
\[
 \mathcal F:\bigl(T^c_{\mathrm{disj}}(sC),b\bigr)
       \longrightarrow
       \bigl(T^c_{\mathrm{disj}}(s\mathcal T),B^{\mathcal T}\bigr),
 \qquad B^{\mathcal T}\mathcal F=\mathcal F b.
\]
The linear term is interval interpolation followed by the compact
current lift $\Gamma_S$.  It induces an injection on cohomology
into the full target.  Its restriction to the polynomial current
subcomplex of the target is an $A_\infty$ quasi-isomorphism.
\end{theorem}
\begin{proof}
The interval recursion can be checked at each finite word length.
Suppose its morphism equation and $b^2=0$ hold below length $n$.
With the length-$n$ terms temporarily zero, write the morphism defect
as $Q_n=R_n-L_n$, where
\[
 L_n=\sum_{\substack{r+q+t=n\\2\le q<n}}
 f_{r+1+t}(1^{\otimes r}\otimes b_q\otimes1^{\otimes t}).
\]
Tensor substitution includes the sign of passing the degree-one
map $b_q$ through the preceding suspended inputs.
Let $O_n$ be the length-$n$ term of $b^2$.
The identity
\[
 B^A(B^Af-fb)+(B^Af-fb)b=-fb^2
\]
gives $\delta Q_n=-IO_n$, with
$\delta Q_n=B^A_1Q_n+Q_n b_1$.
For $j\ge2$, the recursion implies $Pf_j=Hf_j=0$.
Thus $PL_n=HL_n=0$.
Applying $P$ gives $\delta(PR_n)=-O_n$.
Set $b_n=PR_n$ to make $b^2$ vanish in length $n$.
Since $B^A_1H+HB^A_1=1-IP$ and $HI=0$, setting $f_n=-HR_n$
gives
\[
 B^A_1f_n-f_n b_1=-Q_n+IPQ_n+H\delta Q_n=Ib_n-Q_n.
\]
This completes the length-$n$ morphism equation.
Induction proves $B^Af=fb$ and $b^2=0$.
All operations preserve the support labels, so the induction uses
only defined products.

The strict map $G$ induces a coalgebra map commuting with bar
differentials.  Composing it with $f$ proves the displayed identity
for $\mathcal F$ at every length.
The contraction makes $I$ a quasi-isomorphism.
The lift $\Gamma$ is an isomorphism onto the polynomial current
subcomplex, proving the restricted quasi-isomorphism.
For the full target, let $j_S$ extend by zero from $U_S$ to $U$.
Let $\tau_S$ set normal variables outside $E_S$ to zero.
The moment projection $L_U$ of
Proposition~\ref{prop:bf-compact-full-boundary-split} satisfies
$L_U\rho=\pi e^{-\hbar K_Q}$, where
$K_Q=\sum_{\nu<\omega}p_{\nu\omega}\partial_{x_\nu}\partial_{x_\omega}$
and $\pi(x_\nu)=y_\nu$, $\pi(\theta_\nu)=0$, $\pi(y_\nu)=y_\nu$.
This exponential cancels $e^{\hbar K_S}$ on an input using only $E_S$.
Thus $\tau_SL_Uj_S\Gamma_S=1$ and
$p\otimes(\tau_SL_Uj_S)$ is a chain left inverse to $i\otimes\Gamma_S$.
This proves its cohomology injection.
\end{proof}

\section{The first two Taylor coefficients}

Put $v=se_1$, $w=s\eta$, and $\mathbf e=s(e_0+e_1)$, of degrees
$-1,0,-1$, respectively.  For $f_j\in W_{S_j}$ with disjoint labels,
write
\[
 U(f_1,\ldots,f_n)=
 \mu\bigl(\Gamma_{S_1}f_1,\ldots,\Gamma_{S_n}f_n\bigr)
 =\Gamma_{\cup S_j}(f_1\star\cdots\star f_n).
\]
Define
$c(t)=(t-t^2)/2$ and $g(t)=t^3/6-t^2/4+t/12$.
All nonzero binary and ternary Taylor terms on $v,w$ are
\begin{equation}
\label{eq:ordered-current-low-taylor}
 \begin{aligned}
 \mathcal F_2(vf_1,wf_2)&=s\bigl(c(t)U(f_1,f_2)\bigr),\\
 \mathcal F_2(wf_1,vf_2)&=-s\bigl(c(t)U(f_1,f_2)\bigr),\\
 \mathcal F_3(vf_1,wf_2,wf_3)&=s\bigl(g(t)U(f_1,f_2,f_3)\bigr),\\
 \mathcal F_3(wf_1,vf_2,wf_3)&=-2s\bigl(g(t)U(f_1,f_2,f_3)\bigr),\\
 \mathcal F_3(wf_1,wf_2,vf_3)&=s\bigl(g(t)U(f_1,f_2,f_3)\bigr).
 \end{aligned}
\end{equation}
Every other such value is zero.  Values with an interval unit
$\mathbf e$ are zero in arities at least two.  Linearity and
$se_0=\mathbf e-v$ determine every value on $C_I$.
For example, $B^A_2(st,sdt)=s(tdt)$, while
$B^A_2(sdt,st)=-s(tdt)$.  Hence $f_2(v,w)=sc$ and $f_2(w,v)=-sc$.
The three nonzero cubic expressions $R_3$ are
$s(cdt)$, $-2s(cdt)$, and $s(cdt)$.
Since $-h(cdt)=g$, these give \eqref{eq:ordered-current-low-taylor}.
All remaining expressions either vanish or are functions, which
$h$ kills.  A unit input reduces a tree to $HI$, $H^2$, or $PH$,
so its higher Taylor terms vanish.

For homogeneous suspended inputs $u,v,w$, the full first morphism
equations, including the unary input terms, are
\begin{equation}
\label{eq:ordered-current-F2-equation}
 \begin{split}
 B_1^{\mathcal T}\mathcal F_2(u,v)
 +B_2^{\mathcal T}(\mathcal F_1u,\mathcal F_1v)
 ={}&\mathcal F_1b_2(u,v)+\mathcal F_2(b_1u,v)\\
 &+(-1)^{|u|}\mathcal F_2(u,b_1v),
 \end{split}
\end{equation}
and
\begin{equation}
\label{eq:ordered-current-F3-equation}
 \begin{split}
 B_1^{\mathcal T}\mathcal F_3(u,v,w)
 &+B_2^{\mathcal T}(\mathcal F_1u,\mathcal F_2(v,w))
  +B_2^{\mathcal T}(\mathcal F_2(u,v),\mathcal F_1w)\\
 ={}&\mathcal F_1b_3(u,v,w)+\mathcal F_2(b_2(u,v),w)
       +(-1)^{|u|}\mathcal F_2(u,b_2(v,w))\\
 &+\mathcal F_3(b_1u,v,w)+(-1)^{|u|}\mathcal F_3(u,b_1v,w)\\
 &+(-1)^{|u|+|v|}\mathcal F_3(u,v,b_1w).
 \end{split}
\end{equation}
On $(v,w,w)$ the first and third terms on the left are
$s(g' dt\,U)$ and $s(cdt\,U)$.
The right side is $s(dt\,U/12)$.
Thus this equation is exactly $g'(t)+c(t)=1/12$.
On $(v,v,w)$ it is the function identity
$tc=c/2-3g$.  These equations also check the sign contributed by
the first suspended input.

The ordinary degree-$1-n$ Taylor map $\mathsf F_n$ is specified by
\[
 \mathcal F_n(sa_1,\ldots,sa_n)
 =(-1)^{\sum_{j=1}^n(n-j)|a_j|}\,s\mathsf F_n(a_1,\ldots,a_n).
\]
In particular, on singleton currents,
$\mathsf F_2(e_1y_1,\eta y_2)=c(t)u_1u_2$ and
$\mathsf F_3(e_1y_1,\eta y_2,\eta y_3)=-g(t)u_1u_2u_3$.
The negative ternary sign comes from the middle input of degree one.

\section{Comparison with the annular interval inclusion}

Let $\delta_\nu=x_\nu-y_\nu$ and
$K_Q=\sum_{\nu<\omega}p_{\nu\omega}\partial_{x_\nu}\partial_{x_\omega}$.
The compact contraction has
\[
 \Pi=\pi e^{-\hbar K_Q},\qquad
 H_Q=e^{\hbar K_Q}h_0e^{-\hbar K_Q},\qquad
 D_QH_Q+H_QD_Q=1-\iota\Pi,
\]
where $\pi$ sets $\delta_\nu,\theta_\nu$ to zero and
$h_0=q^{-1}\sum_\nu\theta_\nu\partial_{\delta_\nu}$ in positive
total $(\delta,\theta)$-degree $q$, and is zero in degree zero.
These are the maps of Theorem~\ref{thm:bf-compact-quantum-contraction}.
Let $j_S:\mathcal B_\ell(U_S)\to\mathcal B_\ell(U)$ extend by zero.
Define a degree-minus-one map on $W_S$ by
\[
 T_Sf=\rho H_Q\left(\left.e^{\hbar K_S}f\right|_{y=x}\right).
\]
The notation $y=x$ replaces normal symbols by the local variables,
while the annular variables inside $Q$ remain fixed.

\begin{proposition}[The compact linear comparison]
\label{prop:ordered-current-annular-comparison}
For every $f\in W_S$,
\begin{equation}
\label{eq:ordered-current-linear-homotopy}
 D T_Sf=j_S\Gamma_Sf-\iota_U f.
\end{equation}
On the interval tensor product, the map
$\mathcal K_S(\alpha f)=(-1)^{|\alpha|}\alpha\otimes T_Sf$
satisfies
\[
 d_{\mathcal T}\mathcal K_S+\mathcal K_Sd_A
       =(1\otimes j_S)G_S-(1\otimes\iota_U).
\]
At endpoint $a=0,1$, the linear part of the ordered comparison
therefore has a chain homotopy to
$\Phi_a(\alpha f)=\epsilon_a(\alpha)\iota_U(f)$.
\end{proposition}
\begin{proof}
The polynomial $e^{\hbar K_S}f(x)$ is a $D_Q$-cycle.
Applying $\Pi$ cancels its exponential and gives $f(y)$.
The compact contraction proves \eqref{eq:ordered-current-linear-homotopy}.
For the interval identity, the terms containing $d_I\alpha$
have opposite signs and cancel.
The other terms equal $\alpha\otimes DT_Sf$.
Compose with interval inclusion and endpoint evaluation to obtain
the final statement.  Endpoint evaluation kills $dC_I$.
\end{proof}

For two and three distinct current indices, the primitives are
\[
 \begin{aligned}
 T(y_1y_2)&=\rho\left[
  \frac{\theta_1(x_2+y_2)+\theta_2(x_1+y_1)}2\right],\\
 T(y_1y_2y_3)&=\rho\left[
  \sum_{i=1}^3\theta_i\left(
    y_jy_k+\frac{\delta_jy_k+y_j\delta_k}{2}
      +\frac{\delta_j\delta_k}{3}+\frac{\hbar p_{jk}}3
  \right)\right],
 \end{aligned}
\]
where $\{j,k\}=\{1,2,3\}\setminus\{i\}$.
Indeed $H_Qe^{\hbar K_Q}=e^{\hbar K_Q}h_0$ on these inputs.
The second derivative of $\delta_j\delta_k/3$ produces the
positive quantum correction $\hbar p_{jk}/3$.
Their differentials are, respectively,
\[
 \begin{aligned}
 u_1u_2+\hbar p_{12}-e_1e_2,\\
 u_1u_2u_3+\hbar(p_{12}u_3+p_{13}u_2+p_{23}u_1)-e_1e_2e_3.
 \end{aligned}
\]
For the actual product $u_1u_2u_3$, its common-annulus projection is
\[
 e_1e_2e_3-\hbar(p_{12}e_3+p_{13}e_2+p_{23}e_1).
\]
Thus projecting the coefficient of
$\mathsf F_3(e_1y_1,\eta y_2,\eta y_3)$ gives $-g(t)$ times
the complete three-current Wick expression.
Here $e_i$ in the current products denotes the annular form,
whereas $e_1$ in the input denotes the interval endpoint cochain.

The constants $p_{\nu\omega}$ are rational functions of the
distinct insertion points, with every jet derivative retained.
An undivided jet of order $n$ is obtained by multiplying its divided
generator by $n!$.  With $\ell=-k$, the elementary contraction becomes
$\hbar k(a_i-a_j)^{-2}$, as in the free-field interval construction.
At $\ell=0$, all exponential contractions vanish and every displayed
map still exists.

For each fixed finite jet list, all formulas extend coefficientwise
to $\hbar$-adic completions with finite polynomial degree in each
coefficient.  Complete each supported word separately and retain
the direct sum over finite word lengths.
A fixed coefficient receives finitely many terms.
The equations hold modulo every $\hbar^M$, hence in the inverse limit.
No commutation of full-boundary cohomology with completion is required.

\begin{proposition}[The binary equation for annular representatives]
\label{prop:ordered-current-annular-binary-obstruction}
Suppose $\hbar p_{12}\ne0$.  The common-annulus linear map cannot
satisfy the binary ordered morphism equation from the Wick source
if the target binary operation is imposed as ordinary polynomial
multiplication in $\mathcal B_\ell(U)$.
\end{proposition}
\begin{proof}
Use the constant interval cochains $(e_0+e_1)y_1$ and
$(e_0+e_1)y_2$.  Their differentials vanish and their Wick product
is $(e_0+e_1)(y_1y_2-\hbar p_{12})$.
The binary morphism equation would require a degree-minus-one
observable $q$ with
\[
 Dq=\iota_U(y_1y_2-\hbar p_{12})-\iota_U(y_1)\iota_U(y_2)
    =-\hbar p_{12}1.
\]
Applying the chain projection $L_U$ gives $0=-\hbar p_{12}1$
in the polynomial normal algebra, a contradiction.
\end{proof}

The target of Theorem~\ref{thm:ordered-current-taylor} retains the
unions $U_S$ of the original disjoint disks.  Its operations use
their actual factorization products.  Extending all entries to a
common disk and replacing them by annular representatives changes
these support conditions.  Proposition~\ref{prop:ordered-current-annular-comparison}
gives the resulting linear chain homotopy, but it does not turn
overlapping multiplication into an ordered bar differential.
An ordered comparison allowing arbitrary disk changes requires
coherent maps for those changes.  Collision chains require their
own domains and boundary-stratum equations.
